\documentclass{article}
\usepackage{indentfirst,amsmath,graphicx,algorithm,listings}
\usepackage[a4paper, total={6in, 8in}]{geometry}
\usepackage{amssymb}
\usepackage[utf8]{inputenc}
\usepackage{amsmath, amsthm}
\usepackage{xcolor}

\usepackage[english]{babel}
\usepackage{amsthm}
\usepackage[UTF8]{ctex}
\usepackage[colorlinks=true, allcolors=blue]{hyperref}


\newtheorem{theorem}{Theorem}
\newtheorem{definition}{Definition}
\newtheorem{proposition}{Proposition}
\newtheorem{Proof}{Proposition}
\newtheorem{remark}{Remark}
\newtheorem{lemma}{Lemma}
\newtheorem{corollary}{Corollary}
\newtheorem{fact}{Fact}


\title{Convergence and divergence of Fourier series}
\author{Lewei Zhang\\
Department of Mathematics\\
Université Paris-Est Créteil\\
\texttt{}
\and
Supervised by:\\
Mr.Chenmin Sun\\
Department of Mathematics\\
Université Paris-Est Créteil\\
\texttt{}
}






\begin{document}
\clearpage

\maketitle
\tableofcontents

\newpage


\section*{Acknowledgements}
I would like to express my gratitude to Mr.Sun，this is our first attempt at writing a mathematical paper, and while our comprehension of mathematics is still developing，we deeply appreciate his encouragement. 
\par Throughout our weekly meetings, Mr.Sun patiently provided invaluable assistance. His expertise and support played a vital role in our journey.

\newpage


\section{Introduction}
The main goal of this article is to present the paper of Kolmogorov which means the construction of a function  whose partial sums of its fourier series diverge almost everywhere.\par
We start to  review some basic properties of fourier series in section 2, including the definition of partial sum, Dirichlet kernel, convolution, and the relation beside which will be used later.\par
In section 3, we introduce a little about the convergence at the beginning as review too, then we try to construct the $ \mathcal{L} ^1 $ function which mentioned before.\par
Our proof is divided into two parts. The first part involves constructing a sequence of functions that satisfies our expectations conditions, while the second part entails constructing the function we desire.\par
More details can be found in Section 3.2.






\section{Basic properties}
We start from the definition and properties of fourier series.

\begin{definition}%1
    The fourier series is defined by $ \sum\limits_{-\infty}^{\infty} \hat{f}(n)e^{2\pi inx/L}$
\end{definition}

\begin{definition}%2
    The $ N^{th} $ fourier coefficient is given by $ \hat{f}(n) =  \frac{1}{L} \int_{L}f(x)e^{\frac{-2\pi inx}{L}}$
\end{definition}

\begin{remark}%这个可以补充
   We also define the trigonometric polynomial by trigonometric series 
\end{remark}

\begin{definition}%3
    The partial sum $ S_N(f)(x)=\sum\limits_{-N}^{N}\hat{f}(n)e^{2\pi inx/L} $
\end{definition}

\begin{definition}%4
    Dirichlet kernel $ D_N(x)=\sum\limits_{-N}^{N}e^{inx} $.  The kernel functions are periodic with period $ 2\pi $
\end{definition}




\begin{proposition}%1
    We can write  Dirichlet kernel as another form:
    \[  D_N(x)=\sum\limits_{-N}^{N}e^{inx}=\frac{sin(N+\frac{1}{2})x}{\sin(\frac{x}{2})}\]
\end{proposition}
\begin{proof}
    we can separate $ \sum\limits_{-N}^{N}e^{inx} $ into two parts, one is $ \sum\limits_{-N}^{0}e^{inx} $ and $ \sum\limits_{1}^{N}e^{inx} $
    \[  \sum\limits_{-N}^{0}e^{inx} = \frac{1-\omega^{N+1}}{1-\omega}\]
    and \[  \sum\limits_{1}^{N}e^{inx}=\frac{\omega^{-N}-1}{1-\omega} \]
    with $ \omega = e^{ix} $
    so we got \[ D_N(x)= \frac{1-\omega^{N+1}}{1-\omega} +\frac{\omega^{-N}-1}{1-\omega}= \frac{\omega^{\frac{1}{2}}(\omega^{-N-\frac{1}{2}}-\omega^{N-\frac{1}{2}})}{\omega^{\frac{1}{2}}(\omega^{-\frac{1}{2}}-\omega^{-\frac{1}{2}})}   =  \frac{\sin(N+\frac{1}{2})x}{\sin(\frac{x}{2})} \]
\end{proof}

\begin{definition}%5
    Given two $ 2\pi -$periodic integrable functions $ f $ and $ g $, we define their convolution $ f*g $ on $ [-\pi,\pi] $ by 
    \[ f*g(x)=\frac{1}{2\pi} \int_{-\pi}^{\pi} f(y)g(x-y)dy\]
    Also, we can change variables:
    \[ f*g(x)=\frac{1}{2\pi} \int_{-\pi}^{\pi} f(x-y)g(y)dy\]
\end{definition}


\textcolor{olive}{Now we can compose definition 3,4 and 5 :}



\begin{proposition}\label{proposition}
    The partial sum of the fourier series of $ f $ can be written as the convolution of $ f $ and Dirichlet kernel $ D_N $
   \begin{align*}
    S_N(f)(x)&=\sum\limits_{-N}^{N}\hat{f}(n)e^{inx}\\
    &=(f*D_N)(x)
   \end{align*}
\end{proposition}

\begin{proof}
    \begin{align*}
        S_N(f)(x)&=\sum\limits_{-N}^{N}\hat{f}(n)e^{inx}\\
        &=\sum_{-N}^{N}\left(\frac{1}{2\pi}\int_{-\pi}^{\pi}f(y)e^{-iny}dy\right)e^{inx}\\
        &=\frac{1}{2\pi}\int_{-\pi}^{\pi}f(y) \left(\sum_{-N}^{N}e^{in(x-y)} \right)dy\\
        &=(f*D_N)(x)
       \end{align*}
\end{proof}

This proposition will be used later.
\begin{theorem}
    \textbf{Parseval's theorem} states that if a function $f(x)$ is square-integrable over a finite interval $[-L, L]$, meaning that 
    \[
    \int_{-L}^{L} |f(x)|^2 \, dx < \infty 
    \]
    then its Fourier series representation converges in mean square to $f(x)$, and the identity holds:
    \[
    \frac{1}{2L} \int_{-L}^{L} |f(x)|^2 \, dx = \sum_{n=-\infty}^{\infty} |c_n|^2
    \]
    Here, $c_n$ are the Fourier coefficients of $f(x)$, given by:
    \[
    c_n = \frac{1}{2L} \int_{-L}^{L} f(x) e^{-i\frac{n\pi}{L}x} \, dx
    \]
    
\end{theorem}
This theorem essentially establishes a relationship between the energy of a function in the time (or spatial) domain and its frequency domain representation. It's a fundamental result in Fourier analysis.
%一个小问题，是特拉斯函数是L1函数吗\\
%这里定义了good kernel，然后我们写一个定理，说明函数与这个好核的卷积的极限等于函数，我们又知道函数的傅立叶级数的部分和可以写成函数与它的dirichl kernel 的卷积，如果Dn是个好核，那么部分和就收敛于f，但是Dn不是好核。






\newpage

\section{The Convergence of Fourier Series}

\subsection{A little about Convergence }
\begin{definition}%6
    A family of good kernels $ {K_n(x)}_{n=1}^{\infty} $ is said to be \textbf{good kernels} if it satisfies :
    \begin{enumerate}
        \item For all $ n \geqslant 1  $, \[ \frac{1}{2\pi}\int_{-\pi}^{\pi} K_n(x)=1\]
        \item There exist $ M \geqslant 0 $ such that for all $ n \geqslant 1 $,\[\int_{-\pi}^{\pi}\left\lvert K_N(x) \right\rvert  dx \leqslant M\]
        \item For every $ \delta > 0 $ \[ \int_{\delta \leqslant \left\lvert x \right\rvert \leqslant \pi} \left\lvert K_N(x) \right\rvert  dx \rightarrow 0\]
    \end{enumerate}
\end{definition}
\begin{theorem}
    Let $ {K_n(x)}_{n=1}^{\infty}  $ a family of good kernels, f be an integrable function on the circle, then whenever f is continuous at x,\[ \lim_{n\rightarrow \infty}(f*K_N)(x)=f(x)\]
    If f is continuous, then convergence is uniform on $ [-\pi,\pi] $
\end{theorem}

\begin{proof}
    If $ \varepsilon >0 $ and f is continuous at x, choose $ \delta  $  such that if $|y| < \delta $, then $|f(x − y) − f(x)| < \varepsilon $.\\
    Using the first property of good kernels, we can write
    \begin{align*}
        (f*K_n)(x)-f(x)&=\frac{1}{2\pi}\int_{-\pi}^{\pi}K_n(y)f(x-y)dy-f(x)\\
        &=\frac{1}{2\pi}\int_{-\pi}^{\pi}K_n(y)[f(x-y)-f(x)]dy
    \end{align*}

    Taking the absolute value

    \begin{align*}
        \left\lvert(f*K_n)(x)-f(x)\right\rvert&=\frac{1}{2\pi}\left\lvert \int_{-\pi}^{\pi}K_n(y)[f(x-y)-f(x)]dy \right\rvert\\
        &\leqslant \int_{|y|<\delta}\left\lvert K_n(y) \right\rvert |f(x-y)-f(x)|dy\\
        &+\frac{1}{2\pi}\int_{\delta \leqslant |y|\leqslant \pi}\left\lvert K_n(y) \right\rvert |f(x-y)-f(x)|dy\\
        &\leqslant \frac{\epsilon }{2\pi}\int_{-\pi}^{\pi}\left\lvert K_n(y) \right\rvert dy+\frac{2B}{2\pi}\int_{\delta \leqslant |y|\leqslant \pi}\left\lvert K_n(y) \right\rvert dy
    \end{align*}

    where B is a bound for $ |f| $. The second property of good kernels shows that

    \[ \frac{\epsilon }{2\pi}\int_{-\pi}^{\pi}\left\lvert K_n(y) \right\rvert dy \leqslant \frac{\epsilon M}{2\pi}\]
    for all $ n \geqslant 1 $. The third property of good kernels shows that
    \[ \frac{2B}{2\pi}\int_{\delta \leqslant |y|\leqslant \pi}\left\lvert K_n(y) \right\rvert dy \leqslant \epsilon  \]
    for all $ n \geqslant N(\delta ) $. Therefore we have
    \[ | (f*K_n)(x)-f(x)| \leqslant C_{\epsilon }\]





    
\end{proof}

If the Dirichlet kernels were a good family of kernels,then we could develop a sense of convergence from the partial sums of the Fourier series. However, we will see that the family of Dirichlet kernels is not a family of good kernels. In particular, the Dirichlet family of kernels does not satisfy the second property of good kernels.

\begin{proof}
    We can derive this using the geometric series identity and trigonometric identities, \[D_N(x)=\sum\limits_{-N}^{N}e^{inx}=\frac{sin(N+\frac{1}{2})x}{sin(\frac{x}{2})}\]
    Now
     \begin{align*}
        \int_{-\pi}^{\pi}\left\lvert D_N(x) \right\rvert dx &\geqslant c \int_{-\pi}^{\pi}\frac{\left\lvert sin(N+\frac{1}{2})x \right\rvert}{\left\lvert x \right\rvert}\\
        &\geqslant c \int_{\pi}^{N\pi}\frac{|sin(x)|}{|x|}dx\\
        &= c\sum_{k=1}^{N-1}\int_{k\pi}^{(k+1)\pi}\frac{1}{k}\\
        &\geqslant clog(n)
    \end{align*}
    where $c > 0$. This implies that the Dirichlet kernels are not a family of good kernels.
\end{proof}
\textcolor{olive}{So we need another definition to state convergence.}
\par
\begin{definition}
    We define the Cesàro mean of the Fourier series by 
    \[ \sigma_N(f)(x)=\frac{S_0(f)(x)+S_1(f)(x)+\dots S_{N-1}(f)(x)}{N} \]
\end{definition}

\begin{definition}
    The $ N^{th}  $ Fejér kernel is defined by 
    \[ F_N(x)= \frac{D_0(x)+D_1(x)+\dots D_{N-1}(x)}{N}\]
    Fejér kernel is a good kernel.
\end{definition}

\begin{theorem}
    Fejér's theorem states that the Cesàro means of the Fourier series of a function $f(x) $ converge uniformly to $f(x)$ wherever $f(x)$ is continuous.
\end{theorem}

The proof of this theorem is also very complex, and we will not go into it here.





\subsection{Almost everywhere divergence of Fourier series of integrable functions}
In 1922, at the age of 19 while still an undergraduate student at Moscow State University, Kolmogorov constructed an $ \mathcal{L} ^1 $ function in his paper "Une série de Fourier-Lebesgue divergente presque partout," where the partial sums of its Fourier series diverge almost everywhere. 
\par
His proof was very ingenious. We attempted to understand his proof and fill in some details.
\begin{theorem}
    There exist a $ \mathcal{L} ^1 $ function who's fourier series diverge almost everywhere. 
\end{theorem}
We start to proof this theorem by the construction of a sequence of functions $ \phi _n  \subset  \mathcal{L} ^1[0,2\pi] $ which satisfying the following properties:


\begin{enumerate}
    \item\label{C1}$ \phi_n(x)\geqslant 0 $ and $ \int_{0}^{2\pi}\phi_n(x)dx=2 $
    \item The partial sum of Fourier series of $ \phi_n(x) $ is bounded
    \item\label{C3} For every function $ \phi_n(x) $ there exist $ M_n  \in \mathbb{R}  $, a set $ E_n \subset [0,2\pi] $, and $ q_n \in \mathbb{N}  $ such that 
        \begin{enumerate}
            \item $ \lim\limits_{n\rightarrow +\infty}M_n = +\infty$
            \item $ \lim\limits_{n\rightarrow +\infty}\left\lvert E_n \right\rvert = 2 \pi $
            \item $ \forall x \in E_n, \exists q_n \in \mathbb{N}  $ such that for $ k \leqslant q_n $
            \[ \left\lvert S_k\phi_n(x)\geqslant M_n \right\rvert \]
        \end{enumerate}
\end{enumerate}
We note $ C(1),C(2),C(3)$ these three properties. \par
So we now begin to justify the existence of this sequence of functions.
\subsubsection*{Condition 1 and 2}


We take a increasing sequence with odd natural numbers.
\[ \lambda_1=1 \ \  \lambda_2,\dots,\lambda_n \] 

then we define a sequence of natural numbers $ m_k $ s.t
\[
    \begin{cases}
    m_1=n   \\
    2m_k+1 =\lambda_k(2n+1) 
\end{cases} 
\]

a sequence of real numbers $ A_k $
\[
\begin{cases}
    A_k &=\frac{4\pi k}{2n+1}\quad 1 \leqslant k \leqslant n \\
    A_n&= 2\pi- \frac{2\pi}{2n+1}
\end{cases}
\]

a sequence of interval 

\[ \Delta_k=[A_k-\frac{1}{m_k^2},A_k+\frac{1}{m^2_k}] \]

those interval satisfies :

\begin{enumerate}
    \item $ \Delta_k$ mutually disjoint
    \item $ \Delta_k \subset [0,2\pi] $
\end{enumerate}

Here is the demonstration:\par 
\begin{proof}
    \begin{enumerate}
        \item $ \Delta_k  $ is contained in the ball(which is disjoint) centred at $ A_k $ with radius $ \frac{1}{2}|A_k-A_{k-1}| $
        \item  we first proof that $ \frac{1}{m_k^2}<\frac{2\pi}{2n+1} $ holds because:
        \begin{align*}
            m_k = \frac{\lambda_k(2n+1)}{2} - \frac{1}{2} \Rightarrow m_k^2 &= \frac{[\lambda_k(2n+1)-1]^2}{4} \\
            &> \frac{[\lambda_k(2n+1)-1]^2}{2\pi} \\
            &> \frac{2n+1}{2\pi}
        \end{align*}
        then we proof $ \Delta_k \subset [0,2\pi] $ 

        \begin{enumerate}
            \item $ \Delta_k \subset \left[0,+\infty\right) $
                   \[\frac{1}{m_k^2}\leqslant \frac{2\pi}{2n+1} \Rightarrow A_1-\frac{1}{m_k^2}>\frac{4\pi}{2n+1}-\frac{2\pi}{2n+1}=\frac{2\pi}{2n+1}\overset{n \to \infty}{\longrightarrow}0\]
            \item $\Delta_k \subset [0,2\pi ]$
                    \[ A_n +\frac{1}{m_k^2}=\frac{4 n \pi}{2n+1}+\frac{1}{m_k^2}<\frac{4n \pi}{2n+1}+\frac{2\pi}{2n+1}=2\pi \]
                   
        \end{enumerate}
        \end{enumerate}     
\end{proof}

Then we can define a function:

\begin{align}
\phi_n(x) =
\begin{cases}
    \frac{m_k^2}{n}& \text{if } x\in \Delta_k \quad 1 \leqslant k \leqslant n \\
    0 & \text{if } x \in [0,2\pi]\backslash \cup_{k=1}^n \Delta_k
\end{cases}
\end{align}

we now have 

\begin{enumerate}
    \item $ \phi_n(x)\geqslant 0 $
    \item $ \int_{0}^{2\pi}\frac{m_k^2}{n}dx = \sum_{k=1}^{n}\frac{m_k^2}{n}\frac{2}{m_k^2}=2 $
    \item $\phi_n(x) $ $ 2\pi $ periodic also integrable.
\end{enumerate}

So condition 1 holds.\par
Condition 2 is immediate since $ \phi_n(x) $ is a piecewise constant function.\par
\subsubsection*{Condition 3 }



Now let's define a sequence of intervals $ {\sigma_k}_{2 \leqslant k \leqslant n} $

\[ \sigma_k=[A_{k-1}+\frac{2}{n^2},A_k-\frac{2}{n^2}] \]

we must prove that $ A_{k-1}+\frac{2}{n^2}\leqslant  A_k-\frac{2}{n^2} $


\begin{proof}
    \[ A_{k-1}+\frac{2}{n^2}=\frac{4\pi}{2n+1}(k-1)+\frac{2}{n^2}=\frac{4k\pi}{2n+1}-\frac{4\pi}{2n+1}+\frac{2}{n^2} \]
    \[ <\frac{4k\pi}{2n+1}-\frac{2}{n^2}=A_k-\frac{2}{n^2} \]
\end{proof}

\begin{lemma}
    

 $ \sigma_k $ is strictly between $ \Delta_{k-1} $  and $ \Delta_k $
\end{lemma}
\begin{proof}
    we want to show that 
    \[ A_{k-1}+\frac{1}{m^2_{k-1}}<A_{k-1}+\frac{2}{n^2}<A_k-\frac{2}{n^2} <A_k+\frac{1}{m^2_k}\]
    there are three inequality, we need to proof the first and the last, the last one is trivial cause both $ \frac{2}{n^2} $ and $ \frac{1}{m_k^2} $ are positive. And for the first inequality:
   
    \begin{align*}
    m_k^2 = \frac{[\lambda_k(2n+1)-1]^2}{4} \Rightarrow m_k^2 \geqslant \frac{\lambda_k^2(4n^2)}{4}\Rightarrow m_k^2 \geqslant \frac{n^2}{2}\Rightarrow \frac{1}{m_{k-1}^2} \leqslant \frac{2}{n^2} 
    \end{align*}

\end{proof}

\begin{lemma}
    $ \lim\limits_{n\rightarrow \infty} |\cup_{k=1}^n\sigma_k|= 2\pi  $
\end{lemma}

\begin{proof}
 \[ |\sigma_k | = A_k-A_{k-1}-\frac{4}{n^2}= \frac{4\pi}{2n+1}-\frac{4}{n^2}\]
so their union $|\cup_{k=1}^n\sigma_k| = \frac{4\pi n}{2n+1}-\frac{4}{n}$ almost gives the entire interval $ [0,2\pi] $ as $ n $ grows.

\end{proof}

~\\
\textcolor{olive}{From now on, we try to prove the most difficult part of this paper.}

We first claim two theorem which will be used later:

\begin{theorem}
    \textbf{Mean Value Theorem} If $f(x)$ is continuous on $[a, b]$ and differentiable on $(a, b)$, then there exists $c$ in $(a, b)$ such that $f'(c) = \frac{f(b) - f(a)}{b - a}$.

\end{theorem}
\begin{theorem}
    \textbf{The Riemann-Lebesgue lemma} If $f(x)$ is a bounded and integrable function,then its Fourier transform $F(\xi)$ 
    \[
    \lim_{|\xi| \to \infty} F(\xi) = 0
    \]
    This can be expressed as:

\[
\lim_{|\xi| \to \infty} \int_{-\infty}^{\infty} f(x) e^{-i\xi x} \, dx = 0
\]

    
\end{theorem}














~\\

By proposition\ref{proposition} in section 2, we can get that 


\begin{align}
    S_{m_k}(\phi_n)=\frac{1}{2\pi}\int_{0}^{2\pi}\phi_n(\alpha)\frac{\sin\frac{2m_k+1}{2}(\alpha-x)}{\sin\frac{1}{2}(\alpha-x)}\mathrm{d}\alpha
\end{align}

\textbf{A.The choice of $ m_k $}

It is clear to see that 


\begin{align*}
    S_{m_k}(\phi_n)&=\frac{1}{2\pi}\int_{0}^{2\pi}\phi_n(\alpha)\frac{\sin\frac{2m_k+1}{2}(\alpha-x)}{\sin\frac{1}{2}(\alpha-x)}\mathrm{d}\alpha\\
    &=\frac{1}{2\pi}\sum_{s=1}^{k-1}\int_{\Delta_s}\phi_n(\alpha)\frac{\sin\frac{2m_k+1}{2}(\alpha-x)}{\sin\frac{1}{2}(\alpha-x)}\mathrm{d}\alpha \\
    &+ \frac{1}{2\pi}\sum_{s=k}^{n}\int_{\Delta_s}\phi_n(\alpha)\frac{\sin\frac{2m_k+1}{2}(\alpha-x)}{\sin\frac{1}{2}(\alpha-x)}\mathrm{d}\alpha \quad \text{by (1)}\\
    &=I_1+I_2 
\end{align*}

Then we focus on $ \phi_n^{(s)} $, whose intergal will represent the integral on a segment $ \Delta_s $, we then obtain the equality below:
%\textcolor{red}{这里，或许我多想了。它就是第i 或者说第s个}

\begin{align*}
    \frac{1}{2\pi}\int_{\Delta_s}\phi_n(\alpha)\frac{\sin\frac{2m_k+1}{2}(\alpha-x)}{\sin\frac{1}{2}(\alpha-x)}\mathrm{d}\alpha 
    &=\frac{1}{2\pi}\int_{0}^{2\pi}\phi_n^{(s)}\frac{\sin\frac{2m_k+1}{2}(\alpha-x)}{\sin\frac{1}{2}(\alpha-x)}\mathrm{d}\alpha\\
    &= S_{m_k}(\phi_n^{(s)})
\end{align*}
\par 
Assume we have already defined $ \lambda_n $, thus we define $ m_k $, we will discuss the choose of $ m_k $ later.\par 
After we will estimate $ I_1 $ and $ I_2 $.

Now consider the partial sum $ S_{m_k}(\phi_n^{(s)}) $ where $ x \in \sigma_k $.

\begin{enumerate}

\item $ 1 \leqslant s \leqslant k-1 $\par
        Notice that, for any $ x \in \sigma_k $ and $ \alpha \in \Delta_s $, $ x $ and $ \alpha $ belong to two disjointed intervals,so we have:
\[ \exists \epsilon >0 \text{ s.t } |\alpha-x|\geqslant \epsilon \]
    consequently, \[ \frac{1}{\sin\frac{1}{2}(\alpha-x)} \] is a continuous function of $ \alpha  $, so this function is integrable, by the Riemann-Lebesgue lemma:
\[\lim_{|m_k|\rightarrow \infty}S_{m_k}(\phi_n^{(s)}) = 0 \] 
        \begin{proof}
   $  \phi_n^{(s)}$ is a bounded and integrable function on $ \Delta_s $, so we can get his fourier transform tends to zero by the Riemann-Lebesgue lemma, then the partial sum also tends to zero.
        \end{proof}

So we can choose that $ m_k $ big enough,which imply $ \lambda_k $ big enough grand to make 

        \[ |I_1|<1 \]


\item $ s \geqslant k $
\begin{align*}
    \frac{1}{2\pi}\int_{\Delta_s}\frac{m_s^2}{n}\frac{\sin\frac{2m_k+1}{2}(\alpha-x)}{\sin\frac{1}{2}(\alpha-x)}\mathrm{d}\alpha 
    &=\frac{1}{2\pi}\int_{\Delta_s}\frac{m_s^2}{n}\frac{\sin\frac{2m_k+1}{2}(\alpha-x)}{\sin\frac{1}{2}(\alpha-x)}\mathrm{d}\alpha - \frac{1}{\pi}\int_{\Delta_s}\frac{m_s^2}{n}\frac{\sin\frac{2m_k+1}{2}(A_s-x)}{\sin\frac{1}{2}(A_s-x)}\mathrm{d}\alpha\\
    &+\frac{1}{2\pi}\int_{\Delta_s}\frac{m_s^2}{n}\frac{\sin\frac{2m_k+1}{2}(A_s-x)}{\sin\frac{1}{2}(A_s-x)}\mathrm{d}\alpha\\ 
    &= \frac{1}{2\pi}\int_{\Delta_s}\frac{m_s^2}{n}\left[\frac{\sin\frac{2m_k+1}{2}(\alpha-x)}{\sin\frac{1}{2}(\alpha-x)}   - \frac{\sin\frac{2m_k+1}{2}(A_s-x)}{\sin\frac{1}{2}(A_s-x)} \right]\mathrm{d}\alpha \\
    &+\frac{1}{2\pi}\int_{\Delta_s}\frac{m_s^2}{n}\frac{\sin\frac{2m_k+1}{2}(A_s-x)}{\sin\frac{1}{2}(A_s-x)}\mathrm{d}\alpha\\
    &=J_1+J_2
\end{align*}


We first focus on $ J_1 $:\par
\textcolor{olive}{(The following four remarks will be used later)}
\par

\begin{remark}
    We can find that 
    \[ \frac{2m_k+1}{2}(A_s-A_k)=\frac{4\pi(s-k)}{2n+1}\frac{2m_k+1}{2}=2\pi(s-k)\lambda_k \]
    so that 
    \[ \sin\frac{2m_k+1}{2}(A_s-x)=\sin\left[\frac{2m_k+1}{2}(A_s-x)-\frac{2m_k+1}{2}(A_s-A_k)\right]\]
    \[ =\sin\frac{2m_k+1}{2}(A_k-x) \]
\end{remark}

\begin{remark}
    \[ |D_N'(x)|=|\sum\limits_{-N}^{N}ine^{inx}| \leqslant |\sum_{n=0}^{N}n| \leqslant N^2\]
\end{remark}
\begin{remark}
    For all $ \alpha \in \Delta_s $, we have $ |A_s-\alpha|\leqslant \frac{1}{m_s^2} $ by the definition of $ \Delta_s $
\end{remark}

\begin{remark}
    
    For all $\ x \in \sigma_k$,we have 
    \[ 0<A_s-x<A_s-A_{k-1}=\frac{4\pi}{2n+1}(s-k+1)<\frac{2\pi}{n}(s-k+1) \]
\end{remark}

~\\


 By using the mean value theorem, we look 

\[ \frac{\sin\frac{2m_k+1}{2}(\alpha-x)}{\sin\frac{1}{2}(\alpha-x)}   - \frac{\sin\frac{2m_k+1}{2}(A_s-x)}{\sin\frac{1}{2}(A_s-x)}  \]

as the derivative between $ \alpha -x  $ and $ A_s-x $, also, by the remark 3 and 4

\[ | \frac{\sin\frac{2m_k+1}{2}(\alpha-x)}{\sin\frac{1}{2}(\alpha-x)}   - \frac{\sin\frac{2m_k+1}{2}(A_s-x)}{\sin\frac{1}{2}(A_s-x)}| \leqslant |A_s-\alpha||m_k^2| \leqslant \frac{m_k^2}{m_s^2}\leqslant 1\] 

So now we look back at $ J_1 $,

\[ |J_1|\leqslant \frac{1}{2\pi}\int_{\Delta_s}|\frac{m_s^2}{n}\left(\frac{\sin\frac{2m_k+1}{2}(\alpha-x)}{\sin\frac{1}{2}(\alpha-x)}   - \frac{\sin\frac{2m_k+1}{2}(A_s-x)}{\sin\frac{1}{2}(A_s-x)} \right)| \mathrm{d}\alpha  \leqslant \frac{1}{2\pi}\int_{\Delta_s}\frac{m_s^2}{n}\mathrm{d}\alpha\]
\[\leqslant \frac{1}{2\pi}\frac{m_k^2}{n}\frac{2}{m_k^2}=\frac{\tau}{n}  \quad \text{ with } |\tau|<1\] 

%\textcolor{red}{所以这里可以直接把 $ m_s $换成 $ m_k $}


Then $ J_2 $




\begin{align*}
    J_2= \frac{1}{2\pi}\int_{\Delta_s}\frac{m_s^2}{n}\frac{\sin\frac{2m_k+1}{2}(A_s-x)}{\sin\frac{1}{2}(A_s-x)}\mathrm{d}\alpha&=\frac{1}{2\pi}\frac{\sin\frac{2m_k+1}{2}(A_s-x)}{\sin\frac{1}{2}(A_s-x)}\int_{\Delta_s}\frac{m_s^2}{n}\mathrm{d}\alpha\\
    &=\frac{1}{2\pi}\frac{2}{n}\sin\left(\frac{2m_k+1}{2}(A_k-x) \right)\frac{1}{\sin\frac{1}{2}(A_s-x)}\\
    &= \frac{1}{n\pi}\sin\left(\frac{2m_k+1}{2}(A_k-x) \right)\frac{1}{\sin\frac{1}{2}(A_s-x)}\quad  \text{by remark 2}
\end{align*}

Finally we can look back at $ I_2 $
\begin{align*}
    I_2=\sum_{s=k}^{n}(J_1+J_2)=\frac{1}{n\pi}\sin\left(\frac{2m_k+1}{2}(A_k-x) \right)\sum_{s=k}^{n}\frac{1}{\sin\frac{1}{2}(A_s-x)}+C \quad \quad \text{where |C| < 1}
\end{align*}

by the remark 5 we can find that 
\[ \frac{1}{n\pi}\sum_{s=k}^{n}\frac{1}{\sin\frac{1}{2}(A_s-x)} \geqslant \frac{1}{n\pi}\sum_{s=k}^{n}\frac{n}{\pi(s-k+1)}=\frac{1}{\pi^2}\sum_{s=k}^{n}\frac{1}{s-k+1} =\frac{1}{\pi^2}\sum_{r=1}^{n-k+1}\frac{1}{r} \geqslant \frac{1}{\pi^2}\ln(n-k+2) \] 


so 
\[ I_2 \geqslant \frac{1}{\pi^2}\sin\left(\frac{2m_k+1}{2}(A_k-x) \right)\ln(n-k+2)-2  \quad \quad \forall x \in \sigma_k, \ \ 2 \leqslant k \leqslant n\] 

\end{enumerate}
The partial sum equal to $ I_1 +I_2 = I_1+J_1+J_2$ and we have $ I_1 < 1$, and also $J_1 < 1 $, so $ J_2 $ is the main bound.\par
~\\

\textbf{B.Now we can find $ M_n $, $ q_n $ and the set $ E_n$ }
 


In the previous demonstration, we got the bound:
\[| S_{m_k}(\phi_n)|\geqslant \frac{1}{\pi^2}|\sin(m_k +\frac{1}{2})(A_k-x)|\ln(n-k+2) -2 \]

We take $ k \leqslant n - \sqrt{n} $, then

\[  | S_{m_k}(\phi_n)| \geqslant \frac{1}{2\pi^2}|\sin(m_k +\frac{1}{2})(A_k-x)|\ln n-2\]

We set \[ E_n= \bigcup_{2 \leqslant k \leqslant n-\sqrt{n}}E_n^{(k)} \]

where \[E_n^{(k)}=\left\{ x: x \in \sigma_k \ \ \text{ and } \ |\sin(m_k +\frac{1}{2})(A_k-x)| \geqslant \frac{2\pi^2}{\sqrt{\ln n}}\right\}\]

If we choose $ m_k $ large enough 

\[  mes\left\{x: x \in \sigma_k \ \  |\sin(m_k +\frac{1}{2})(A_k-x)| < \frac{2\pi^2}{\sqrt{\ln n}}\right\}=\mathcal{O}(\frac{1}{n\sqrt{\ln n}})\]

so we have 

\[ mes\  E_n^{(k)} = mes \ \sigma_k- \mathcal{O}(\frac{1}{n\sqrt{n}}) \quad \quad as \ \ n \rightarrow \infty\]

Thus 

\[ mes\ E_n= \sum_{2 \leqslant k \leqslant n- \sqrt{n}}mes\ \sigma_k- \mathcal{O}(\frac{1}{\sqrt{\ln n}})  = (\frac{4\pi}{2n+1}-\frac{4}{n^2})(n-\sqrt{n})-\mathcal{O}(\frac{1}{\sqrt{\ln n}})\]
\[ = 2\pi - \mathbf{o} (1) \]

which imply that 

\[ mes \ E_n=2\pi \ \ \ \ as \  \ n \rightarrow \infty \]
So finally we obtain 
\begin{enumerate}
    \item $ | S_{m_k}(\phi_n)| \geqslant \sqrt{\ln n } -2 = M_n$,  where $ M_n >0 $, $ M_n \rightarrow \infty $ as $ n \rightarrow \infty $
    \item $q_n=m_n$
\end{enumerate}

That is the condition 3 we want.\par
Those three conditions state that there exist a function of bounded variation whose $ \mathcal{L}^1  $ norm is 2 but, at every point of a set whose measure nearly equate to $ 2\pi $, their Fourier sums are arbitrarily large.\par



\subsubsection*{Existence of an $ \mathcal{L}^1  $  function with a.e divergent series}
Now we come to the construction of the $ \mathcal{L}^1 $ function whose Fourier series diverges almost everywhere.

%算了不等他了，我自己写。

$ \phi_n $ has already been constructed, so now we define an increasing sequence of natural numbers $ n_k $ such that:
\begin{enumerate}
    \item $ \frac{1}{\sqrt{M_{n_k}}}<\frac{1}{2^k} $  $\Rightarrow \sum_{k=1}^{\infty}\frac{1}{\sqrt{M_{n_k}}}< 1 $ 
    \item If $ B_n = \sup_{N \geqslant 0}|S_N\phi_n(x)|$, then $ \sum_{i=1}^{k-1}B_{n_i}<\frac{1}{2}\sqrt{M_{n_k}} $
    \item $ 2q_{n_i}+1 \leqslant \frac{\sqrt{M_{n_k}}}{2^k} $ $ \forall i < k  $
\end{enumerate} 
\begin{proof}

We can find a sequence satisfy $ 1,2,3 $ cause 
\begin{enumerate}
   
    \item $ M_n \rightarrow \infty $ and the $ \phi_n $ is bounded variation, $ B_n $ finite.
    \item $ n_s $ are known for all values of $ s  $  less than k.
    
    \item if we have already chosen  $ n_1,\dots ,n_{k-1} $, then by taking $ n_k >n_{k-1} $ sufficiently large, we can satisfy 2 and 3, because $  M_n \rightarrow \infty$ as $ n \rightarrow \infty $
\end{enumerate}
\end{proof}

Now consider the function series $ \sum_{i=1}^{\infty}\frac{1}{\sqrt{M_{n_i}}}\phi_{n_i} (x)$, from the $ C$(\ref{C1}) of function $ \phi_n $
\[ \sum_{i=1}^{\infty}\frac{1}{\sqrt{M_{n_i}}}\int_{0}^{2\pi}\phi_{n_i} (x) = \sum_{i=1}^{\infty}\frac{2}{\sqrt{M_{n_i}}}<2\sum_{i=1}^{\infty}\frac{1}{2^k} < 2\]
from Beppo Levi’s lemma ,this function of series converge almost everywhere on $ [0,2\pi] $ to some function $ \Phi(x) $, so we can define 
\[ \Phi(x)= \sum_{i=1}^{\infty}\frac{1}{\sqrt{M_{n_i}}}\phi_{n_i} (x)\] which will be the function whose Fourier series diverge almost everywhere.\par
The partial sum 
\[ S_N(\Phi(x))=\sum_{-N}^{N}\frac{1}{2\pi}e^{inx}\int_{0}^{2\pi}\Phi(x)e^{-iny}\mathrm{d}y  =\sum_{i=1}^{\infty}\frac{1}{\sqrt{M_{n_i}}}S_N(\phi_{n_i}(x))\]
\par
We choose a natural number $ k $ and a point $ x\in E_{n_k} $, and let us consider the Fourier sums corresponding to $ N=N_k $
\begin{align*}
    |S_{N_k}(\Phi(x))|&=|\sum_{i<k}\frac{1}{\sqrt{M_{n_i}}}S_{N_k}(\phi_{n_i}(x))+\sum_{i>k}\frac{1}{\sqrt{M_{n_i}}}S_{N_k}(\phi_{n_i}(x))+\frac{1}{\sqrt{M_{n_i}}}S_{N_k}(\phi_{n_i}(x))|\\
    &=|L_1+L_2+L_3|\\
    &\geqslant |L_3| -|L_2+L_1|
\end{align*}
 
by using the condition we just defined:
\[ |L_1|= \sum_{i<k}\frac{1}{\sqrt{M_{n_i}}}S_{N_k}(\phi_{n_i}(x)) \leqslant \sum_{i<k}\frac{1}{\sqrt{M_{n_i}}}\sup|S_{N_k}\left(\phi_{n_i}\left(x\right)\right)| \leqslant \sum_{i=1}^{k-1}\frac{B_{n_i}}{\sqrt{M_{n_i}}}\]
\[ <\sum_{i=1}^{k-1}B_{n_i}<\frac{1}{2}\sqrt{M_{n_k}}\]


and 
\[ |L_2|= \sum_{i>k}\frac{1}{\sqrt{M_{n_i}}}|S_{N_k}(\phi_{n_i}(x))|\leqslant \sum_{i>k}\frac{1}{\sqrt{M_{n_i}}}(\frac{1}{2}+q_{n_k})\frac{1}{2\pi}\int_{0}^{2\pi}|\phi_{n_i}(x)|\mathrm{d}x\]
\[ = \sum_{i>k}\frac{1+2q_{n_k}}{2\sqrt{M_{n_i}}}\frac{1}{\pi}\leqslant \frac{1}{\pi}\sum_{i>k}\frac{1}{2^i}<\sum_{i>k}\frac{1}{2^i}<\frac{1}{2^k}\]
\par
also
\[ |L_3|= \frac{1}{\sqrt{M_{n_i}}}S_{N_k}(\phi_{n_i}(x)) \overset{C(\ref{C3})(c)}{\geqslant } \sqrt{M_{n_k}} \]

Finally we get
\begin{align*}
    |S_{N_k}(\Phi(x))| &\geqslant |L_3| -|L_2+L_1|\\
    &\geqslant \sqrt{M_{n_k}}-\frac{1}{2}\sqrt{M_{n_k}}-\frac{1}{2^k}\\
    &\geqslant \frac{1}{2}\sqrt{M_{n_k}}-\frac{1}{2}
\end{align*}
which holds for all $ x\in E_{n_k} $.\par
Now let $ E= \limsup\limits_k E_{n_k}  $，which can also be expressed $ E=\bigcap\limits_{j=1}^{\infty} \bigcup\limits\limits_{k=j}^{\infty}  E_{n_k} $, then:

\[ 2\pi \geqslant \bigcup\limits\limits_{k=j}^{\infty}  |E_{n_k}|  \geqslant \lim\limits_l|E_{n_l}| \overset{C(\ref{C3})(b)}{=} 2\pi \Rightarrow \bigcup\limits\limits_{k=j}^{\infty} | E_{n_k}|=2\pi \ \forall k\]
Because it is a numerable intersection of sets with measure $ 2\pi $, all of them contained in $ [0,2\pi] $\par
Then 
\[ |E| = |\bigcap\limits_{j=1}^{\infty} \bigcup\limits\limits_{k=j}^{\infty}  E_{n_k}|=2\pi\]


Finally, we get that for all $ x \in E $ 
\[  |S_{N_k}(\Phi(x))| \geqslant\frac{1}{2}\sqrt{M_{n_k}}-\frac{1}{2} \]
and $ M_{n_k} \rightarrow \infty $, then $ \limsup\limits_k |S_{N_k}(\Phi(x))| = \infty$.\par
So we finally complete our proof: the partial sums of its Fourier series diverge almost everywhere.





%我们定义的是阶梯函数，区间有n个。

%思路重复一下，我们把部分和拆分成两部分，一部分是I1，另一部分是I2，他们分别对应着n个区间中大于等于k 和小于k的部分，然后我们来观察其中的一项，再求和。
























\begin{thebibliography}{9}
\bibitem{citation_key}
Andrey Kolmogorov .(1923)“Une śerie de Fourier-Lebesgue divergente presque partout”. \textit{Fundamenta Mathematicae}, \textbf{Volume 4 }(Issue 1), Page 324-328.
\bibitem{citation_key}
Elias M. Stein.Fourier analysis an introduction.Publisher :Princeton University Press 2003.
\bibitem{citation_key}
P. L. Ul’yanov. (1983). Kolmogorov and divergent Fourier series. \textit{Russian Mathematical Surveys}, \textbf{Volume83}(Issue4, 57–100), DOI: 10.1070/RM1983v038n04ABEH004205
\bibitem{citation_key}
Tooryanand Seetohul (2023). Convergence of Fourier series and Kolmogorov’s theorem. 

\end{thebibliography}

\end{document}
